\chapter{NDNSF 框架}
\label{chap:ndnsf-ch}

\section{设计目标}

NDNSF 的设计目标如下：

\begin{enumerate}[leftmargin=*]
  \item \textbf{通用服务抽象。}应用应能用统一 service name、请求对象和
  响应对象表达服务，而不是依赖为每个服务生成的静态 stub。
  \item \textbf{数据中心安全。}服务消息和权限对象应作为 NDN Data 被签名、
  加密和验证。
  \item \textbf{协作与韧性。}框架应支持多个提供者、选择性 ACK、自定义选
  择和补偿请求。
  \item \textbf{低延迟路径。}对已知提供者的控制命令，应通过 Targeted
  invocation 减少消息轮次，同时保持权限和 token 验证。
  \item \textbf{操作验证。}框架应支持 MiniNDN 评估、选择性真机实验、证
  书 bootstrap 和可复现配置。
\end{enumerate}

\section{NDNSF v0.1 基线}

NDNSF v0.1 包括：

\begin{itemize}
  \item ServiceController、ServiceUser 和 ServiceProvider 运行时；
  \item 通用动态 C++ 服务 API；
  \item 基于 NDN Data 和 Sync 的 Request、ACK、Selection 和 Response
  消息；
  \item 基于 NAC-ABE 的服务级访问控制；
  \item 提供者权限和用户权限获取；
  \item 选择性 ACK 和自定义选择策略；
  \item 自适应 admission control；
  \item 与 NSC 和 gRPC 的 MiniNDN 对比评估。
\end{itemize}

这个 baseline 的意义在于，它已经把数据中心服务框架必须协调的主要部分连
接起来：service naming、service-message publication、用户和提供者授权、
加密服务 payload、provider metadata、provider selection 和 response
verification。换言之，NDNSF v0.1 不只是一个 API 草图，而是一个可运行的
runtime，展示了如何把 NDN name、Sync、signed Data、NAC-ABE attribute 和
service-level policy 组合成一个完整服务调用 workflow。因此，本 proposal
可以把重点放在扩展和验证 framework，而不是从零开始设计。

图~\ref{fig:ndnsf-v01-architecture-ch} 总结了 v0.1 baseline 中的三个运
行时角色。Controller 管理 policy 并分发权限；Provider 订阅有权限处理的
请求并发布服务结果；User 发布请求、收集 ACK metadata、选择 provider 并
验证 response。这个角色划分是后续 UAV service container 和分布式推理服
务协作的基础。

\begin{figure}[H]
\centering
\begin{subfigure}[t]{0.78\textwidth}
  \centering
  \includegraphics[width=\linewidth]{figures/NDNSF_ServiceController.png}
  \caption{Controller}
\end{subfigure}

\vspace{0.8em}

\begin{subfigure}[t]{0.78\textwidth}
  \centering
  \includegraphics[width=\linewidth]{figures/NDNSF_ServiceProvider.png}
  \caption{Provider}
\end{subfigure}

\vspace{0.8em}

\begin{subfigure}[t]{0.78\textwidth}
  \centering
  \includegraphics[width=\linewidth]{figures/NDNSF_ServiceUser.png}
  \caption{User}
\end{subfigure}
\caption{作为 proposal baseline 的 NDNSF v0.1 运行时角色。}
\label{fig:ndnsf-v01-architecture-ch}
\end{figure}

\section{作为核心扩展的服务提供者协作}

NDNSF v0.1 之后最核心的扩展是 provider collaboration。在协作服务中，一个
逻辑 request 不一定只由单个 endpoint 完成，而可以由多个已授权 provider
共同参与。Provider 可以发布 ACK metadata，例如可用性、队列状态、位置或
role capability；User 可以在 ACK collection interval 之后选择一个或多个
provider；当一个逻辑任务只完成了一部分时，应用可以发出 compensation
request。对于耗时较长的 selection，User 还应能诊断被选中的 provider 当
前是未收到 selection、排队中、执行中、出错，还是已经发布 response。

这个协作模型应保持 application-independent。多无人机巡逻可以用它分配任
务段并从部分失败中恢复；分布式推理可以用它分配模型 role 并协调中间
artifact。两类应用都不应把自己的领域语义写入 NDNSF core。NDNSF 提供的
通用语义是 request publication、provider metadata、provider selection、
authorization、status diagnosis 和 response verification；具体领域工作由
应用定义。

\section{通用动态 API}

NDNSF 当前使用通用动态 API。服务提供者注册 typed handler：

\begin{lstlisting}[language=C++]
provider.addHandler<RequestT, ResponseT>(
    ndn::Name("/ObjectDetection/YOLOv8"),
    handler);
\end{lstlisting}

服务用户请求服务：

\begin{lstlisting}[language=C++]
user.RequestService<RequestT, ResponseT>(
    providers,
    ndn::Name("/ObjectDetection/YOLOv8"),
    request,
    onResponse,
    onTimeout,
    timeoutMs,
    strategy);
\end{lstlisting}

这个 API 将应用对象与框架消息分离。应用数据位于 typed request/response
对象或 RequestMessage/ResponseMessage payload 中。框架负责命名、发布、
加密、验证、provider selection 和 timeout。

\section{服务消息模型}

NDNSF v0.1 使用四类框架消息，将服务发布、provider admission、provider
selection 和结果返回分开。这一消息模型是 one-to-many service design 的
核心：同一个请求可以被多个有权限的 provider 观察到，provider 可以在真正
执行之前暴露可用性和运行时 metadata，用户则显式选择一个或多个 provider。

\begin{table}[H]
\centering
\caption{NDNSF v0.1 服务消息。}
\label{tab:ndnsf-message-model-ch}
\begin{tabularx}{\textwidth}{p{0.30\textwidth}p{0.16\textwidth}X}
\toprule
消息 & 生产者 & 作用 \\
\midrule
RequestMessage & User & 发布命名服务请求和应用 payload。 \\
RequestAckMessage & Provider & 表示 provider 是否愿意处理请求，并携带用于
选择的运行时 metadata。 \\
SelectionMessage & User & 选择一个或多个 provider，并用 provider token
将请求绑定到被选中的 provider。 \\
ResponseMessage & Provider & 被选中的 provider 返回服务执行结果。 \\
\bottomrule
\end{tabularx}
\end{table}

\begin{table}[H]
\centering
\caption{四类消息携带的代表性字段。}
\label{tab:ndnsf-message-fields-ch}
\begin{tabularx}{\textwidth}{p{0.30\textwidth}X}
\toprule
消息 & 代表性字段 \\
\midrule
RequestMessage & service name、request ID、requester identity、payload、
UserToken \\
RequestAckMessage & service name、request ID、provider identity、status、
message、metadata payload、回显的 UserToken、ProviderToken \\
SelectionMessage & service name、request ID、selected provider identity、
ProviderToken \\
ResponseMessage & service name、request ID、provider identity、result
payload、回显的 UserToken \\
\bottomrule
\end{tabularx}
\end{table}

这些字段保持 service-independent。应用相关数据留在 payload 或 typed
request/response 对象中；框架则使用 service name、request ID、identity
和一次性 token 完成授权检查、防重放和 response verification。

\section{服务调用流程}

普通 NDNSF 服务调用包括以下步骤：

\begin{enumerate}[leftmargin=*]
  \item 用户发布 RequestMessage，并携带一次性 UserToken。
  \item 有权限的提供者解密请求，检查 provider permission，并根据 ACK 策
  略决定是否回应。
  \item 提供者发布 RequestAckMessage，返回 UserToken、一次性
  ProviderToken 和 ACK metadata。
  \item 用户在 ACK timeout 内收集候选者，根据策略选择一个或多个提供者。
  \item 用户发布 SelectionMessage，并携带被选中提供者的 ProviderToken。
  \item 提供者验证 ProviderToken，执行服务并发布 ResponseMessage。
  \item 用户验证 response 中的 UserToken 并交付应用结果。
\end{enumerate}

该流程支持一对多发现、应用自定义选择和防重放保护。对已知提供者的低延迟
调用，Targeted API 先通过 normal bootstrap 获取一批 token pair，然后后续
请求可走 REQUEST 到 RESPONSE 的 fast path。

\section{安全模型}

NDNSF 使用 controller 授权的权限模型。ServiceController 维护用户和提供者
policy，向目标 identity 生成加密的 PermissionResponse。服务消息根据服务
级授权 namespace 映射到 NAC-ABE attribute。一次性 user/provider token 将
调用流程绑定到特定 request，防止 replay。

提交的原型将这些一次性字段实现为 \texttt{UserToken} 和
\texttt{ProviderToken}；修订后的协议描述使用 \texttt{RequestNonce} 和
\texttt{ProviderChallenge} 表达其请求绑定和选择绑定作用。二者都不是
NDN Interest Nonce，也不是通用的持有者凭证。

图~\ref{fig:ndnsf-policy-example-ch} 展示了 NDNSF 中 provider policy 和
user policy 的关系：provider permission 决定一个实体可以提供哪些命名服
务，user permission 决定一个实体可以访问哪些命名服务。本文后续 UAV 和
DistributedInference 扩展都沿用这种以 service name 为中心的 policy 模型，
同时补充更明确的 trust root、publisher namespace 和部署流程。

\begin{figure}[H]
\centering
\includegraphics[width=0.82\textwidth]{figures/policy-examples.png}
\caption{NDNSF provider 和 user access-control policy 示例。}
\label{fig:ndnsf-policy-example-ch}
\end{figure}

该模型区分：

\begin{itemize}
  \item \textbf{用户授权}：用户是否可以请求某服务；
  \item \textbf{提供者授权}：提供者是否可以提供某服务；
  \item \textbf{消息机密性}：只有具有相应 attribute 的实体可以解密服务
  消息；
  \item \textbf{消息完整性和来源认证}：NDN Data signature 和 trust schema
  保护消息；
  \item \textbf{防重放}：UserToken 和 ProviderToken 一次性使用。
\end{itemize}

\subsection{采用 ABE-backed 服务授权的理由}
\label{sec:abe-rationale-ch}

NDNSF 采用 ABE-backed 授权，是为了支持面向授权服务提供者的保密发现，
并将服务权限管理与身份签名凭证分离。这是针对 NDNSF 调用流程的设计取舍，
并非 ABE 全面优于角色授权的结论。执行授权决定实体是否可以调用或提供某项
服务；内容机密性决定谁能读取受保护的消息内容。签名验证、服务权限检查和
调用状态检查仍是不同的要求。

\paragraph{发现阶段的保密性。}
在选择 Provider 之前，User 发布 discovery descriptor（发现描述），其中
包含 Provider 决定是否返回 positive ACK 所需的位置、截止时间和资源要求。
例如，附近 UAV 需要知道观察区域和到达时限才能决定是否参与，但此时不需要
完整应用输入。ABE 无需逐一列出 Provider 身份，即可将发现描述的读取权限
限制为获准提供该服务的 Providers。这里的保密对象是缺少相应解密能力的外部
观察者；获准但最终未被选中的 Providers 仍可读取。内容加密不隐藏外层 Data
名称或流量模式。在请求级机密性路径中，完整应用输入及其解密材料通过
Selection 限定给选中的
Providers。Trust Schema 本身不加密内容；角色授权方案也可以配合 ABE 或
服务组加密来实现发现阶段的保密性。

\paragraph{权限聚合与凭证复用。}
KP-ABE 将属性关联到密文，将访问策略编码到解密密钥
（DKEY）中~\cite{dulal2023nacabe}。NDNSF 使用
\texttt{/SERVICE/<service>} 表示 Provider 的服务解密权限，使用
\texttt{/PERMISSION/<service>} 表示 User 的调用相关解密权限。在同一授权
机构和 ABE 参数代际内，一个实体的多项服务权限被合并为 OR 策略，编码到
单独为该实体签发的 DKEY 中，同时复用其身份签名凭证。DNMP 则通过角色签名
密钥和 Trust Schema 约束命令的目标与动作~\cite{dnmp}。已有角色可以覆盖多项
权限；不同权限需求可能要求调整角色分配、凭证或规则，但不意味着每项服务
都必须建立新角色。

作为对照，我们提出的每服务角色证书方案会为实体获准承担的每项服务角色签发
证书。User 使用相应服务角色密钥签名 Request 和 Selection Data，Provider
使用相应密钥签名 ACK 和 Response Data。验证方检查服务、消息类型、证书
所授角色、证书状态及调用状态，内容保密则由配套加密机制提供。ABE 避免了
这些额外的服务专用签名凭证，但 DKEY 仍需签发和更新，其大小及处理成本可能
随策略复杂度增长。证书方案也能复用凭证，因此凭证对象数量较少本身不足以
证明总成本更低。

\paragraph{为何继续使用 ABE-backed challenges？}
Requester Challenge（实现中的 \texttt{UserToken}）和 Provider Challenge
（\texttt{ProviderToken}）用于检查参与者在当前交换中是否具有相应的服务和
调用权限解密能力。签名标识参与身份，匹配检查和一次性状态检查拒绝不匹配或
重放的消息。这沿用了属性化权限管理，而不必另外增加服务角色签名凭证，但
并非没有额外成本：保密发现已需要 Provider DKEY，而 User 的调用权限 DKEY
和 challenge 处理仍是额外授权成本。ABE 保密配合角色证书授权同样是有效的
替代方案。

\subsection{运行时权限撤销}
\label{sec:authorization-withdrawal-ch}

当前运行时通过 Controller 签名的授权状态以及调用处理点上的检查实现权限
撤销。参与者需要获取并验证更新后的状态，才能执行新公布的权限撤销。
ABE 参数轮换和 DKEY 更新保护随后使用新代际加密的交换。目前轮换作用于整个
Controller 的 ABE 参数代际，其他仍获授权的成员也可能需要刷新 DKEY；
这不是无需更新其他成员的撤销方案。

DNMP 描述了通过轮换加密密钥限制后续测量数据访问，但没有详细给出可对应
比较的在线命令授权撤销协议~\cite{dnmp}。角色证书对照方案可以撤销相关凭证，
并要求执行节点检查签名的撤销状态或拒绝过期证书。这是在验证方取消凭证
所代表的权限，而不是物理取回已分发的证书。两类方案都需要传播更新信息，
也都无法收回已披露的密钥或明文。

定向集成测试已检验 NDNSF 的运行时撤销检查，包括身份级、证书级和服务级
撤销。这些测试支持已覆盖运行时边界上的行为，不代表整个网络中的瞬时撤销，
也不能证明其成本低于证书方案。后续评估应测量状态传播延迟、拒绝生效延迟、
重新分发密钥的流量，以及对未撤销成员的影响。因此，运行时状态验证和密钥
轮换使 NDNSF 支持权限撤销，同时带来明确的状态传播和重新分发密钥的成本。


\section{已完成基础：NDNSF v0.1 评估}

已完成的 NDNSF 评估在 MiniNDN 中将 NDNSF 与 NSC 和 gRPC 进行对比。评估
关注有线、丢包、间歇提供者、无线移动、自适应 admission、选择性 ACK 和
自定义选择场景中的服务延迟和成功率。证据支持的结论是有条件的：在丢包、
间歇 provider 和单一 endpoint 绑定容易失效的场景中，provider 多样性可以
提高可靠性；这并不等价于 NDNSF 在所有动态环境中都具有更高的终端成功率或
更低的延迟。

最有说服力的证据来自 endpoint binding 或单一 provider 变脆弱的场景。在
Memphis--CSU--UCLA 的有线丢包路径中，如果每条 link 的 one-way loss 是
$\ell$，packet round-trip loss 是 $1-(1-\ell)^4$。1\% link loss 时，
packet round-trip loss 为 3.94\%，NDNSF 完成 599/600 个 request；3\%
link loss 时，packet round-trip loss 为 11.47\%，NDNSF 在启用 SVS retry
后仍完成 598/600 个 request。10\% link loss 时，packet round-trip loss
上升到 34.39\%，最佳 NDNSF 结果为 539/600。这说明高丢包下的主要边界是
SVS-based service invocation 的 bounded-time delivery，而不是 NDNSF service
model 本身失败。在有线多 provider 的间歇失效实验中，provider failure probability 为 20\% 时，
NDNSF 仍保持 99.46\% success rate；失效概率为 30\% 时仍有 96.70\%，而单
provider baseline 会下降到约三分之二或更低。无线移动表中的数值属于已投稿
NDNSF v0.1 的历史诊断，拓扑、provider 数量和计时来源与当前公平的 Spec171
四-provider campaign 不完全一致，因此不能把它们作为当前 mobility superiority
的直接证据。当前校正后的单 AP、50\,m 覆盖、2\,m/s 移动、三 seed 复现实验中，
NDNSF success rate 为 75.33\%，相对 gRPC 的 71.78\% 提高 3.55 个百分点，
相对 NSC 的 70.11\% 提高 5.22 个百分点；但在 50\,m/15\,m/s 和
100\,m/15\,m/s 条件下并不占优，因此这里应表述为条件性结果，而非普遍
mobility 优势。Adaptive
admission 和 Selective ACK 进一步说明 runtime metadata 的价值：过载时
admission control 将 P95 latency 保持在 270\,ms 以下；当最近 provider 过
载时，Selective ACK 将 success rate 从 1.82\% 提高到 100\%。这些结果共同
支撑较窄的结论：one-to-many invocation、provider metadata 和显式 selection
能在保持数据中心安全的同时，为特定 provider-diversity 和间歇可用性场景提供
可靠性收益；当前证据还不足以支持普遍的无线 mobility 优势。

\begin{table}[H]
\centering
\caption{从已投稿 NDNSF v0.1 历史评估中总结的代表性可靠性结果；无线行不作为当前公平 mobility superiority 的证据。}
\label{tab:ndnsf-v01-reliability-ch}
\begin{tabularx}{\textwidth}{p{0.27\textwidth}p{0.21\textwidth}ccc}
\toprule
场景 & 条件 & NSC & gRPC & NDNSF \\
\midrule
有线多 provider 间歇失效 & 10\% 失效概率 & 83.25\% & 83.35\% & \textbf{99.82\%} \\
有线多 provider 间歇失效 & 20\% 失效概率 & 66.58\% & 66.69\% & \textbf{99.46\%} \\
有线多 provider 间歇失效 & 30\% 失效概率 & 49.92\% & 65.36\% & \textbf{96.70\%} \\
无线移动 provider & 100\,m AP range & 45.42\% & 45.23\% & \textbf{96.25\%} \\
无线移动 provider & 150\,m AP range & 68.75\% & 68.05\% & \textbf{100.00\%} \\
无线移动 provider & 200\,m AP range & 88.33\% & 87.14\% & \textbf{100.00\%} \\
\bottomrule
\end{tabularx}
\end{table}

\begin{table}[H]
\centering
\caption{从已投稿 NDNSF v0.1 评估中总结的代表性机制结果。}
\label{tab:ndnsf-v01-mechanisms-ch}
\begin{tabularx}{\textwidth}{p{0.24\textwidth}p{0.29\textwidth}X}
\toprule
机制 & Baseline 行为 & NDNSF 行为 \\
\midrule
100 RPS offered load 下的 adaptive admission & 关闭 admission control 时，P95 latency 达到 4254.39\,ms，timeout rate 达到 3.55\%。 & 开启 admission control 后，实际 admitted rate 限制在约 81 RPS，P95 latency 为 268.12\,ms，timeout rate 为 0.00\%。 \\
最近 provider 过载时的 Selective ACK & ACK-all 机制把 2250 个请求中的 2209 个选到最近但过载的 provider，success rate 只有 1.82\%。 & Selective ACK 抑制过载 provider 的 ACK，把 selection 转移到可用 provider，并达到 100.00\% success rate。 \\
真机间歇可用性实验 & 单 provider NSC/gRPC 在绑定 provider 半程不可用时 success rate 接近 50\%。 & NDNSF 在两个 provider 中动态选择，10 RPS 下达到 74.62\%，30 RPS 下达到 73.73\%，接近预期可用窗口。 \\
\bottomrule
\end{tabularx}
\end{table}

\section{NDNSF v0.1 的局限}

NDNSF v0.1 证明了方向可行，但还没有完全覆盖受控框架测试之外的动态服务
应用需求。
推动后续博士论文工作的主要局限包括：

\begin{itemize}
  \item \textbf{已知提供者的低延迟控制}：normal
  Request/ACK/Selection/Response workflow 适合 provider discovery，但发
  往已知无人机或设备的控制命令需要更快的 Targeted path，同时仍要保留授
  权和防重放保护。
  \item \textbf{服务提供者协作}：v0.1 支持 provider selection，但更大的
  任务需要多个 provider 协作、处理部分失败并发出 compensation request。
  \item \textbf{Service-container 组合}：真实应用通常在同一受信任进程内
  包含多个服务。这些服务需要高效本地组合，但不能向远程 caller 暴露绕过
  网络安全机制的路径。
  \item \textbf{大 artifact}：v0.1 service payload 不适合承载模型、录制
  视频、图片或中间张量。应用应在 service payload 中放置 typed object
  reference，再使用 NDN 原生 segmented retrieval 获取大型签名和加密对象。
  \item \textbf{操作验证}：v0.1 展示了安全模型，但更广泛的验证还需要更
  清晰的证书 bootstrap、trust-schema 指导、配置流程和可复现实验。
  \item \textbf{应用级验证}：v0.1 的 MiniNDN 评估证明核心框架可运行，但
  博士论文还需要评估该设计是否能支撑 UAV 和分布式推理这类高压力 workflow。
\end{itemize}

\section{Provider Collaboration 的支撑机制}

以下机制用于支撑 provider collaboration 在特定运行条件下成立。它们由 UAV
和分布式推理 workload 驱动，但不是独立博士论文主线，也不应成为应用专用
shortcut，而应保持为通用 framework mechanism。

\subsection{Targeted Invocation}

对于面向已知 provider 的控制服务，例如 MAVLink 命令，完整 ACK/Selection
流程会增加延迟。Targeted invocation 使用 normal 流程进行首次 bootstrap，
获取一批一次性 token pair，随后每个 request 携带未使用的 ProviderToken，
response 返回对应 UserToken。这样可以减少消息轮次，同时保留权限验证和防
重放语义。

\subsection{Trusted Local Invocation}

Service container 内部的服务组合不需要走网络。同一受信任进程内的一个服
务组件调用另一个内部服务组件时，不应发布 NDNSF Request。拟实现的本地调
用机制支持同进程 trusted local invocation，包括同步和异步调用。该路径不
生成 wire message，不使用 NAC-ABE，也不消耗 token；外部节点不能要求
provider 使用该路径，它只用于同一受信任进程内的服务组合。

\subsection{大数据引用抽象}

服务请求 payload 不应承载模型、视频、图片或大张量。拟采用的通用模式是：
数据生产者先把对象分段发布为具名 NDN Data，然后在 NDNSF request 或
response 中只放 typed object reference。该 reference 可以包含对象名字、
对象类型、version 或 session id、digest、expected size、final segment 信
息、source type 以及可选加密 metadata。接收方收到 request/response 后，
使用 NDN 原生 segmented fetching，例如 ndn-cxx SegmentFetcher 逻辑，按名
字拉取对象。

该抽象应作为 NDNSF 的通用机制，而不是某个特定存储服务的 API。小型
request/response payload 仍适合控制消息和紧凑 typed object；当 payload
变成模型、视频、activation 或日志这类大 artifact 时，同一个逻辑服务调用
应携带 reference，而不是直接携带字节。应用或 helper 可以把 payload 转换
成命名 Data segment，对这些 segment 签名并可选加密，然后通过 NDNSF 传递
typed reference。这样，服务选择、授权、token 检查和 timeout 诊断仍属于
普通 NDNSF 路径，而大字节传输则走 NDN 原生 named-data retrieval 路径。

同一种 reference 格式应能覆盖 UAV 录制、相机 artifact、模型分片、runtime
artifact、中间 activation 和日志。因此，UAV-APP 与 DistributedInference
可以共享同一套框架级大数据词汇：object type、object name、producer、
session 或 version、size、digest、final segment hint 和可选 encryption
metadata。这避免在 service payload 中嵌入应用特定的裸 Data-name 字符串，
也让不熟悉 NDN packet 细节的应用能够审查大对象行为。

该设计可以选择使用持久化存储，但存储在本 proposal 中只是大对象实现支撑
组件，而不是单独的研究主线。如果使用持久化存储，它应保存 opaque 的 NDN
Data segment；这些 Data 已由应用或框架 helper 命名、签名，并可选加密。
安全性不由存储专用接口单独承担，而是由 NDNSF 通用混合加密、签名、token、
权限和 trust-schema 验证保护 service message 及其引用的数据。Discovery
metadata 仍可作为应用层发现和验证加速工具存在，但大对象传输本身应基于
typed object reference 和 NDN 原生分段获取。

本章通过预览可复用 NDNSF 机制来支持 RQ2，包括安全服务调用、provider 授
权、provider collaboration 以及相关支撑机制。它也支持 RQ3，因为 provider
collaboration 会引入可靠性、安全性和延迟 tradeoff，而支撑机制需要在不削
弱核心服务语义的前提下控制开销。
